\documentclass[11pt]{article}
\usepackage[a4paper,margin=28mm]{geometry}
\usepackage{amsmath,amssymb,amsthm,mathtools,hyperref}
\newcommand{\ddc}{dd^c}
\newcommand{\MA}{\operatorname{MA}}
\newcommand{\PSH}{\operatorname{PSH}}
\newcommand{\PP}{\mathbb P}
\newtheorem{theorem}{Theorem}
\newtheorem{lemma}[theorem]{Lemma}
\newtheorem{proposition}[theorem]{Proposition}
\theoremstyle{remark}\newtheorem{remark}[theorem]{Remark}
\title{Density-root spaces and exact $L^1$ equations\newline
on products with a closed factor}
\author{}\date{}
\begin{document}\maketitle
\noindent\textbf{Scope.} The closed-factor results below have undergone
independent reconstruction and mathematical comparison. Neither
unrestricted question of S04 is resolved by this article.
\begin{abstract}
Let a compact Hermitian manifold be multiplied by a normalized compact
K\"ahler manifold. Every nonzero nonnegative integrable density from the closed
factor has an exactly computed positive individual threshold and an
actual finite-mass equation in the pullback of the closed factor's
full-mass class. The original density-root space contains a closed
linear copy of the factor's $L^N$, with an exact norm and uniformly
vanishing root tails. These conclusions apply when the product has
zero lower mass. An explicit projective example lies outside every
higher $L^r$ space and has logarithmic tail decay. The results concern
factor densities, rather than arbitrary densities on the product.
\end{abstract}

\section{Background conventions and product data}
Use $\ddc=i\partial\bar\partial/(2\pi)$ and ordinary wedge powers.
All quasi-psh functions have their genuine canonical upper
representatives. The extended measure is the non-pluripolar product
with smooth-background corrections. For a positive Hermitian form $h$
and a smooth positive probability $dV$, write
\[
 v_-(h)=\inf_{u\in\PSH(X,h)\cap L^\infty}\int_X(h+\ddc u)^n,
 \qquad b_h(f)=\lim_{j\to\infty}c_h(\min(f,j)).
\]
For nonzero nonnegative bounded measurable data on a positive Hermitian
background, continuous bounded equations and unique positive scalars
$c_h$ follow from \cite[Theorem 0.1]{KN} and
\cite[Corollary 0.2]{Ng}. Their $\ddc$ is twice ours, so their
potentials are multiplied by two; a smooth positive reference-volume
ratio is absorbed into the density. They are not applied to unbounded
$L^1$ data. We use the supplied maximal-subsolution characterization
of $b_h$ and positive-threshold attainment on the product, whose
dimension is at least three. Closed-factor existence in every $k\geq2$
is independently supplied by \cite[Theorem A, Theorem 4.6]{GZ}.
It is decreasing in the density and satisfies
$b_h(tf)=b_h(f)/t$, $t>0$. For bounded data, $b_h=c_h$.
These statements are background assumptions, not new existence
theorems established here.

For real or complex measurable roots set
\[
 \rho_h(p)=b_h(|p|^n)^{-1/n},\quad\rho_h(0)=0,
 \qquad
 \mathcal D_h=\{p\in L^n(dV):\rho_h(p)<\infty\}.
\]
The original full test formula is
\[
 \rho_h(p)=\sup_T\int_X |p|k_T\,dV,
 \qquad k_T=\left(\frac{\gamma^n}{dV}\right)^{(n-1)/n},
\]
where the supremum uses \emph{all} smooth positive tests
$T=\gamma^{n-1}$ with $\ddc T=0$ and $\int h\wedge T=1$.
The norm is $\|p\|_{L^n}+\rho_h(p)$, and $\mathcal D_h^0$ is the
closure of bounded measurable roots in this norm.

Let $(X_0,h_0)$ be connected compact Hermitian of dimension $m$, and
let $(Y,\kappa)$ be connected compact K\"ahler of dimension $k\geq2$,
normalized by $\int_Y\kappa^k=1$. Put
\[
 N=m+k,\quad C=\binom Nm,\quad V_0=\int_{X_0}h_0^m,
 \quad c_0=CV_0,\quad H=h_0+\kappa.
\]
All factor forms and functions are pulled back when appropriate. On
$X=X_0\times Y$ put
\[
 \mu=\frac{h_0^m}{V_0}\wedge\kappa^k,
 \qquad R=\frac\mu{dV},\qquad 0<a=\min R\leq R\leq A=\max R,
\]
where $dV$ is any smooth positive probability on $X$. Define
\[
 \mathcal E(Y,\kappa)=
 \left\{v\in\PSH(Y,\kappa):
       \int_Y\MA_\kappa^{\rm ext}(v)=1\right\}.
\]
This definition refers to actual non-pluripolar mass, not an
envelope-equality replacement.

\begin{theorem}\label{th:main}
For every nonzero nonnegative $g\in L^1(Y,\kappa^k)$, with
$M_g=\int_Yg\kappa^k$, one has
\[
 b_H(Rg)=\frac{c_0}{M_g}>0.
\]
There is a genuine $u=\pi_Y^*v$, with
$v\in\mathcal E(Y,\kappa)$ and $\sup v=0$, such that
\[
 \MA_H^{\rm ext}(u)=b_H(Rg)Rg\,dV,
 \qquad\int_X\MA_H^{\rm ext}(u)=c_0.
\]
The linear map $J\xi=R^{1/N}\pi_Y^*\xi$ satisfies, for every real
or complex $\xi\in L^N(Y,\kappa^k)$,
\[
 \|J\xi\|_{L^N(dV)}=\|\xi\|_{L^N(\kappa^k)},\qquad
 \rho_H(J\xi)=c_0^{-1/N}\|\xi\|_{L^N(\kappa^k)}.
\]
Its image is a closed linear subspace of $\mathcal D_H^0$.
If $v_-(h_0)=0$, then $v_-(H)=0$ as well.
\end{theorem}

The theorem computes the original threshold and norm, with all tests.
It does not prove $\mathcal D_H^0=\mathcal D_H$ or the exact equation
for roots outside the displayed factor subspace.

\section{Bounded normalizers and the original truncations}
Let $g\geq0$ be bounded measurable and nonzero. The
bounded-density equation above on the positive closed background $\kappa$
has a continuous bounded genuine potential $v$ satisfying
\[
 \MA_\kappa(v)=c_\kappa(g)g\kappa^k.
\]
Closed bounded Stokes gives $\int\MA_\kappa(v)=1$, hence
$c_\kappa(g)=1/M_g$. Bounded local products and factor dimensions give
\[
 \MA_H(\pi_Y^*v)=C h_0^m\wedge\MA_\kappa(v)
                  =\frac{c_0}{M_g}Rg\,dV.
\]
Uniqueness of the bounded normalization scalar therefore proves
\begin{equation}\label{eq:bounded}
 c_H(Rg)=\frac{c_0}{M_g}.
\end{equation}
Zeros of $g$ are retained; no positive minimum is used.

For arbitrary nonzero $g\geq0$ in $L^1$, the exact product clips obey
\[
 R\min(g,j/A)\leq\min(Rg,j)\leq R\min(g,j/a).
\]
The bounded identity and scalar monotonicity imply
\[
 \frac{c_0}{\int\min(g,j/a)\kappa^k}
 \leq c_H(\min(Rg,j))
 \leq\frac{c_0}{\int\min(g,j/A)\kappa^k}.
\]
Both bounds tend to $c_0/M_g$. Applying closed bounded Stokes to the
clips on $Y$ likewise gives $b_\kappa(g)=1/M_g$.

\section{A mass upper bound makes the factor subsolution exact}
The probability $(g/M_g)\kappa^k$ is absolutely continuous with respect
to smooth volume, and pluripolar sets have smooth volume zero. Thus it
charges no pluripolar set. By \cite[Theorem A, Theorem 4.6]{GZ},
there is a canonical $\kappa$-psh $v$ in the actual full-mass class
solving the factor equation, in every $k\geq2$. Their $dd^c$ is twice
ours: if their potential is $\varphi$, set $v=2\varphi$. The
non-pluripolar definitions agree under the cofinal real-level
truncations $\max(v,-j)=2\max(\varphi,-j/2)$.
In particular a genuine subsolution exists with
\[
 \MA_\kappa^{\rm ext}(v)\geq\frac{g}{M_g}\kappa^k.
\]
Adding a constant makes $\sup v=0$. For $v_j=\max(v,-j)$, closed
bounded Stokes gives $\int_Y\MA_\kappa(v_j)=1$. Consequently
\[
 \int_Y\mathbf1_{\{v>-j\}}\MA_\kappa(v_j)\leq1,
 \qquad\int_Y\MA_\kappa^{\rm ext}(v)\leq1.
\]
The second inequality follows from the defining increasing
non-pluripolar truncation limit. The prescribed lower measure already
has total mass one. The measure
\[
 \MA_\kappa^{\rm ext}(v)-(g/M_g)\kappa^k
\]
is nonnegative with zero mass, so
\begin{equation}\label{eq:factor}
 \MA_\kappa^{\rm ext}(v)=\frac{g}{M_g}\kappa^k,
 \qquad v\in\mathcal E(Y,\kappa).
\end{equation}
This uses the closed factor's mass upper bound. An arbitrary
Hermitian attained subsolution can have excess or infinite mass.

Set $u=\pi_Y^*v$. Locally it is psh plus smooth with its canonical
values and $H+\ddc u\geq0$. Its exact truncations are
$\max(u,-j)=\pi_Y^*v_j$. Bounded local products give
\[
 \MA_H(\max(u,-j))=C h_0^m\wedge\MA_\kappa(v_j).
\]
The level indicator is the pullback of $\mathbf1_{\{v>-j\}}$.
Multiplication and the defining non-pluripolar limit therefore yield
the actual identity
\[
 \MA_H^{\rm ext}(u)=C h_0^m\wedge\MA_\kappa^{\rm ext}(v)
                   =\frac{c_0}{M_g}Rg\,dV.
\]
Its mass is $c_0$. No capacity convergence is inferred from scalar or
measure convergence. The rigorously specified solution class here is
$\pi_Y^*\mathcal E(Y,\kappa)$; no general smaller-background or
stretching assertion is deduced merely from its definition.

\section{The full norm and its uniformly vanishing tails}
For $\xi\in L^N(Y)$ put $g=|\xi|^N$. Since
$|J\xi|^N=Rg$, integration against $dV$ gives the stated $L^N$ norm
identity. The exact threshold proved above gives the stated
$\rho_H$ identity, including zero by convention. Linearity of $J$
applies the same calculation to differences:
\[
 \|J\xi-J\eta\|_{\mathcal D_H}
 =(1+c_0^{-1/N})\|\xi-\eta\|_{L^N(Y)}.
\]
Bounded approximations of $\xi$ in $L^N(Y)$ have bounded images,
so $J\xi\in\mathcal D_H^0$. Completeness of $L^N(Y)$ and the norm
identity make this image closed in $\mathcal D_H$ and hence in
$\mathcal D_H^0$.

The actual amplitude tail satisfies
\[
 |J\xi|\mathbf1_{\{|J\xi|>K\}}
 \leq\bigl(Rg\mathbf1_{\{g>K^N/A\}}\bigr)^{1/N}.
\]
The original full-test formula is monotone in absolute value.
Applying the exact factor threshold to its right side gives
\begin{equation}\label{eq:tail}
 \rho_H(J\xi\mathbf1_{\{|J\xi|>K\}})
 \leq\left(\frac1{c_0}
      \int_{g>K^N/A}g\kappa^k\right)^{1/N}\longrightarrow0.
\end{equation}
A zero tail has norm zero, with no division by a zero density mass.
Its $L^N(dV)$ tail also tends to zero by integrability. If $dV=\mu$,
then $R=1$ and the exact formula is
\[
 \rho_H(\xi\mathbf1_{\{|\xi|>K\}})
 =\left(\frac1{c_0}\int_{|\xi|>K}|\xi|^N\kappa^k\right)^{1/N}.
\]

If $v_-(h_0)=0$, take bounded $h_0$-psh $w_j$ with masses tending
to zero. The bounded local product expansion from this factor gives
\[
 \int_X\MA_H(\pi_{X_0}^*w_j)
       =C\int_{X_0}\MA_{h_0}(w_j)\longrightarrow0.
\]
Thus $v_-(H)=0$, completing Theorem~\ref{th:main}.

\section{A projective endpoint beyond higher integrability}
Let $Y=\PP^k$ with the normalized Fubini--Study form $\kappa$
in the displayed homogeneous coordinates, so its local potential
on $Z_0=1$ is $\log(1+\sum_{i=1}^k|z_i|^2)$.
For $0<\alpha<1$ put
\[
 \psi=\log\frac{\sum_{i=1}^k|Z_i|^2}{\sum_{i=0}^k|Z_i|^2},
 \qquad v_\alpha=1-(1-\psi)^\alpha,
 \qquad\MA_\kappa^{\rm ext}(v_\alpha)=q_\alpha\kappa^k.
\]
The increasing convex composition has derivative at most $\alpha$,
so $v_\alpha$ is genuine $\kappa$-psh. Its only pole is
$o=[1:0:\cdots:0]$. In the affine chart $Z_0=1$, write
$r=\sum_{i=1}^k|z_i|^2$ and
\[
 F(r)=\log(1+r)+1-
       [1-\log(r/(1+r))]^\alpha,\quad
 Q(r)=rF'(r)=\frac{r+\alpha[1-\log(r/(1+r))]^{\alpha-1}}{1+r}.
\]
Radial Stokes with the normalized Fubini--Study constant gives
$\int_{a<|z|^2<b}(\ddc F)^k=Q(b)^k-Q(a)^k$.
Here $Q(0)=0$ and $Q(\infty)=1$; the infinity hyperplane is smooth
and has zero measure. Thus $\int_Yq_\alpha\kappa^k=1$.
Near the pole, with $L=1-\log r$, the tangential and radial
eigenvalues are comparable to $r^{-1}L^{\alpha-1}$ and
$r^{-1}L^{\alpha-2}$. Therefore
\[
 q_\alpha(r)\asymp r^{-k}L^{-1-k(1-\alpha)}.
\]
Radial integration gives $q_\alpha\in L^1$; for every $p>1$ it
leaves an exponentially increasing factor $e^{k(p-1)L}$, so
$q_\alpha\notin L^p$. With $dV=\mu$, the root
$p_\alpha=q_\alpha^{1/N}$ belongs to $\mathcal D_H^0$ and $L^N$,
but to no $L^r$, $r>N$.

Away from a fixed small pole neighborhood $q_\alpha$ is bounded.
Inside it, $\log q_\alpha=kL-[1+k(1-\alpha)]\log L+O(1)$.
For large $K$ the density level set $\{q_\alpha>K^N\}$ is
sandwiched between radial sets $\{L>c_2\log K\}$ and
$\{L>c_1\log K\}$ for fixed $0<c_1<c_2$.
Integrating $L^{-1-k(1-\alpha)}$ and using the exact tail formula gives
\[
 \rho_H(p_\alpha\mathbf1_{\{p_\alpha>K\}})
       \asymp(\log K)^{-k(1-\alpha)/N}.
\]

A concrete zero-lower-mass factor is the elliptic threefold
\[
 X_0=\{[z:w]\in\PP^3:z,w\ne0\}/\langle[z:w]\mapsto[z:e^{2\pi}w]\rangle,
 \qquad z,w\in\mathbb C^2.
\]
For $h_0=\alpha_0+\beta_0+dt\wedge d^ct$, with
$t=\log(\|w\|^2/\|z\|^2)$ and normalized base forms, the supplied
prior construction \cite[Appendix]{Product} proves
$v_-(h_0)=0$ and $V_0=24\pi$. Its complete factor proof is prior
work, rather than a new assertion established in this article.
For its product with $\PP^2$, $c_0=240\pi$ and every probability
factor density has this positive threshold, its uniformly vanishing
root tails and an exact finite-mass equation as above.

\begin{remark}
The new conclusions concern every integrable closed-factor density,
including zeros and arbitrary smooth reference probabilities, and the
entire indicated linear root subspace. They do not exhibit a root
outside $\mathcal D_H^0$, prove equality of the full spaces, or prove
that root-tail control alone supplies the general zero-lower-mass
exact equation. No uniqueness of unbounded potentials or literature
priority is claimed.
\end{remark}
\begin{thebibliography}{9}
\bibitem{KN} S. Ko\l odziej and N. C. Nguyen,
\emph{Weak solutions to the complex Monge--Amp\`ere equation on
Hermitian manifolds}, \href{https://arxiv.org/abs/1312.5491v2}{arXiv:1312.5491v2}, Theorem 0.1.
\bibitem{Ng} N. C. Nguyen,
\emph{The complex Monge--Amp\`ere type equation on compact Hermitian
manifolds and applications},
\href{https://arxiv.org/abs/1501.00891v1}{arXiv:1501.00891v1}, Corollary 0.2.
\bibitem{GZ} V. Guedj and A. Zeriahi,
\emph{The weighted Monge--Amp\`ere energy of quasiplurisubharmonic
functions}, \href{https://arxiv.org/abs/math/0612630v2}{arXiv:math/0612630v2},
Theorem A, Definition 1.1 and Theorem 4.6.
\bibitem{Product} \emph{Exact equations for $L^1$ densities on products
with zero lower mass}, checked standalone source (2026), Appendix;
\href{https://github.com/MingchenXia/non-kahler-pluripotential-theory/blob/7a244bc4938c5b918f7d906c7201f9385e4188f4/research/papers/P07-zero-lower-mass-product-equations.tex}{immutable source and complete prior factor proof}.
\end{thebibliography}
\end{document}
